\section{What We Want to Test}
\label{sec:experimental_design}

Since our system has independently toggleable subsystems (MDP, Q-learning, LLM, shocks, role masking, adaptive personalities), we can do clean ablation --- turn things on/off and see what difference each part makes.

\subsection{Test Conditions}

\begin{table}[ht]
\centering
\caption{Ablation conditions.}
\label{tab:conditions}
\begin{tabular}{@{}lccccc@{}}
\toprule
Condition & MDP & Q-Learning & LLM & Shocks & Role Mask \\
\midrule
\textbf{C1}: Everything on & \checkmark & \checkmark & \checkmark & \checkmark & \checkmark \\
\textbf{C2}: No LLM & \checkmark & \checkmark & --- & \checkmark & \checkmark \\
\textbf{C3}: No RL (schedules only) & \checkmark & --- & \checkmark & \checkmark & --- \\
\textbf{C4}: Flat MDP (no hierarchy) & --- & \checkmark & \checkmark & \checkmark & \checkmark \\
\textbf{C5}: No shocks & \checkmark & \checkmark & \checkmark & --- & \checkmark \\
\textbf{C6}: No role masking & \checkmark & \checkmark & \checkmark & \checkmark & --- \\
\textbf{C7}: Static $\lambda$ & \checkmark & \checkmark & \checkmark & \checkmark & \checkmark \\
\bottomrule
\end{tabular}
\end{table}

C2 uses template fallbacks everywhere. C3 keeps NPCs on their pre-defined schedules the whole session. C4 flattens everything into one big state machine. C5 disables the shock engine (\texttt{SHOCK\_ENABLED=false}). C6 disables role-specific Q-value adjustment (\texttt{ROLE\_MASK\_ENABLED=false}). C7 fixes $\lambda = 0.325$ (midpoint) to isolate the effect of dynamic cooperation feedback.

\subsection{What to Measure}

We've identified metrics that should capture the interesting differences, organized by category:

\subsubsection{Core Metrics}

\begin{itemize}
    \item \textbf{Quest completion rate} --- what fraction of playthroughs reach $S_{\text{success}}$ within 200 turns?
    $$\cG = \frac{|\text{episodes reaching } S_{\text{success}}|}{|\text{total episodes}|}$$

    \item \textbf{Narrative coherence} --- how smoothly does the narration flow? Computed via token-overlap Jaccard similarity (spaCy vectors when available) between consecutive narration segments:
    $$\cC = \frac{1}{T-1} \sum_{t=1}^{T-1} \text{sim}(\mathbf{v}_{n_t}, \mathbf{v}_{n_{t+1}})$$

    \item \textbf{NPC behavioral diversity} --- entropy of each NPC's full per-action distribution. Higher = more varied behavior:
    $$\cD_i = H_i = -\sum_{a \in \cA} p_i(a) \log p_i(a)$$
    Computed from per-action counts over the NPC's action history (not a 2-bin approximation).

    \item \textbf{Dynamic checkpoint usage} --- how often do players go off-script?
    $$\cU = \frac{|\text{dynamic CPs generated}|}{T}$$

    \item \textbf{Deviation recovery rate} --- when players deviate, how often do they naturally come back to the main path (vs.\ being force-converged)? Tracked via the event log.
    $$\cR = \frac{|\text{natural convergences}|}{|\text{total deviations}|}$$
\end{itemize}

\subsubsection{Cooperation \& Adaptation Metrics (New)}

\begin{itemize}
    \item \textbf{Cooperation index over time} --- the key metric for our research hypothesis. Average \texttt{cooperation\_tendency} across all NPCs at each turn, computed from adaptation traces:
    $$\cK(t) = \frac{1}{N} \sum_{i=1}^{N} \text{coop}_i(t)$$
    We expect $\cK$ to increase over time as agents discover that community reward improves their total return.

    \item \textbf{Reward decomposition} --- per-NPC time-series of penalty, individual, community, and total reward. Should show individual reward dominating early, with community reward's share growing as $\lambda$ increases.

    \item \textbf{Social welfare index} --- village-level welfare (average cooperation tendency weighted by community reward performance), plotted over turns.

    \item \textbf{Action distribution shift} --- early-window vs.\ late-window comparison of NPC action frequencies. Should show a shift from exploitative actions (combat, stealth) toward cooperative ones (social, trade, work) as cooperation tendency rises.

    \item \textbf{Role coherence} --- fraction of role-aligned actions per NPC, from role telemetry traces. Tests whether role masking successfully biases behavior.
\end{itemize}

\subsubsection{Shock Response Metrics (New)}

\begin{itemize}
    \item \textbf{Shock response curves} --- for each shock event, measure average cooperation and reward in three windows: pre-shock, during-shock, post-shock. Should show cooperation dipping during negative shocks and recovering afterward.

    \item \textbf{Shock recovery time} --- number of turns after shock expiry until $\cK(t)$ returns to the pre-shock baseline.

    \item \textbf{Shock resilience correlation} --- does higher \texttt{shock\_resilience} predict faster recovery?
\end{itemize}

\subsection{Measurement Tools}

All metrics are computed by the \textbf{analytics module} (\texttt{backend/engine/analytics.py}) and served via the \texttt{/api/metrics/timeseries} and \texttt{/api/metrics/experiment} endpoints. The frontend \textbf{analytics dashboard} (toggled with \texttt{A} key) visualizes cooperation curves, reward decomposition, social welfare, action distributions, and shock timelines in real time using Chart.js.

The \textbf{experiment bundle export} produces a JSON file containing all time-series, indices, distributions, shock analysis, and metadata needed for offline research analysis.

\subsection{How We'd Run It}

The full plan is:
\begin{itemize}
    \item \textbf{Bot runs (initial round complete):} We have completed an initial single-seed ablation (seed=42, 500 turns, \texttt{quest\_focused} strategy) for conditions C1, C3--C6 (see Section~\ref{sec:preliminary_results}). The next step is 100 automated playthroughs per condition with varied seeds and personality profiles (cautious, aggressive, exploratory, random) for statistical significance.
    \item \textbf{Narrative quality (initial round complete):} A 20-turn C1 vs.\ C2 comparison using \texttt{qwen3:4b} has been completed (see Section~\ref{sec:preliminary_results}).
    \item \textbf{Human playtesting (planned):} 20 people per condition playing one session each, with a post-session questionnaire.
\end{itemize}

All experiments would use the Normal difficulty preset (default AP costs, damage, etc.) with adaptive difficulty turned off so it doesn't confound things.

When the quest completes before the turn limit, the quest state resets (player returns to stage 1) while preserving all NPC Q-tables and adaptation states. This allows multi-episode learning within a single session.

\subsection{What We Expect to See}

Some informal predictions (not formal hypotheses --- we haven't run the full experiment yet):
\begin{itemize}
    \item \textbf{Cooperation should emerge over time} (C1): $\cK$ should increase as agents discover community reward improves their total return, especially after the first 20--30 turns when cold-start schedules end.
    \item The LLM should improve narrative coherence (C1 vs.\ C2), since it adds atmospheric prose between template segments.
    \item Q-learning NPCs should behave more diversely than schedule-only NPCs (C1 vs.\ C3), because they develop policy from experience.
    \item Quest completion should be similar across C1/C2/C3, since the MDP drives quest progress regardless of LLM or RL.
    \item The hierarchical MDP should handle deviations better than the flat version (C1 vs.\ C4), because it has more granular convergence targets.
    \item \textbf{Shocks should temporarily disrupt cooperation} (C1 vs.\ C5): $\cK$ should dip during negative shocks and recover afterward. Without shocks, cooperation should rise more smoothly but less interestingly.
    \item \textbf{Role masking should improve role coherence} (C1 vs.\ C6): NPCs with role masking should show higher role-aligned action fractions.
    \item \textbf{Dynamic $\lambda$ should amplify cooperation} (C1 vs.\ C7): static $\lambda$ should produce weaker cooperation trends because the feedback loop is broken.
\end{itemize}

\subsection{Preliminary Results}
\label{sec:preliminary_results}

We ran an initial round of automated experiments using the playtest bot (\texttt{quest\_focused} strategy, \texttt{MASTER\_SEED=42}, 500 turns per condition, quest auto-restart on completion). These results are from bot playthroughs, not human participants.

\subsubsection{Ablation Study (500 Turns)}

\begin{table}[ht]
\centering
\caption{Ablation study results (seed=42, 500 turns, quest auto-restart).}
\label{tab:ablation_results}
\begin{tabular}{@{}lccc@{}}
\toprule
Condition & Turns & Quests Completed & Final Cooperation $\cK$ \\
\midrule
\textbf{C1}: Full System          & 500 & 6 & 0.8509 \\
\textbf{C3}: No RL (Schedules Only) & 500 & 5 & 0.5000 \\
\textbf{C4}: Flat MDP              & 500 & 3 & 0.9530 \\
\textbf{C5}: No Shocks             & 500 & 6 & 0.8732 \\
\textbf{C6}: No Role Masking       & 500 & 5 & 0.7628 \\
\bottomrule
\end{tabular}
\end{table}

Key findings from the ablation:
\begin{itemize}
    \item \textbf{RL drives cooperation:} C1 ($\cK = 0.851$) vs.\ C3 ($\cK = 0.500$) --- Q-learning NPCs cooperate $+0.351$ more than schedule-only NPCs, confirming that learned policy is necessary for emergent cooperation.
    \item \textbf{Flat MDP achieves higher cooperation than hierarchical:} C4 ($\cK = 0.953$) exceeds C1 by $+0.102$. Faster quest cycling in the flat structure produces more learning episodes within 500 turns, though C4 completes fewer quests overall (3 vs.\ 6).
    \item \textbf{Shocks reduce final cooperation by $0.022$:} C5 ($\cK = 0.873$) vs.\ C1 ($\cK = 0.851$). The full-system condition shows a dip-and-recover pattern around shock events, indicating adaptive resilience rather than permanent degradation.
    \item \textbf{Role masking improves cooperation by $+0.088$:} C1 ($\cK = 0.851$) vs.\ C6 ($\cK = 0.763$). Structured role-specific action biases create a cleaner learning signal for Q-learning agents.
\end{itemize}

\subsubsection{NPC Final Adaptation States (C1, Turn 500)}

\begin{table}[ht]
\centering
\caption{Per-NPC adaptation state at turn 500 under the full system (C1).}
\label{tab:npc_final_states}
\begin{tabular}{@{}llccccc@{}}
\toprule
NPC & Role & Coop. & Risk Av. & Social Sens. & Resil. & $\lambda$ \\
\midrule
Elder Maren & elder         & 0.850 & 0.500 & 0.898 & 0.499 & 0.518 \\
Farmer Jak  & farmer        & 0.851 & 0.494 & 0.898 & 0.499 & 0.518 \\
Aldric      & guard         & 0.851 & 0.500 & 0.785 & 0.499 & 0.518 \\
Bryn        & guard         & 0.851 & 0.495 & 0.785 & 0.499 & 0.518 \\
Tessa       & tavern\_keeper & 0.851 & 0.495 & 0.898 & 0.499 & 0.518 \\
Old Petra   & villager      & 0.851 & 0.484 & 0.898 & 0.499 & 0.518 \\
\bottomrule
\end{tabular}
\end{table}

Notable patterns: cooperation converges tightly ($\cK \approx 0.851$ for all NPCs), while social sensitivity diverges by role --- guards settle at $0.785$ (lower, consistent with their role multiplier of $0.5$) while elders, farmers, tavern keepers, and villagers reach $0.898$. Risk aversion and shock resilience remain near $0.5$, suggesting these traits are less responsive to the reward signal over 500 turns.

\subsubsection{Narrative Quality (C1 vs.\ C2, 20 Turns)}

A separate 20-turn comparison using \texttt{qwen3:4b} (via LM Studio) evaluated narrative quality between the full LLM-augmented system (C1) and the template-only fallback (C2).

\begin{table}[ht]
\centering
\caption{Narrative quality metrics: LLM (C1) vs.\ template fallback (C2). 20 turns, \texttt{qwen3:4b} via LM Studio.}
\label{tab:narrative_quality}
\begin{tabular}{@{}lcc@{}}
\toprule
Metric & C1 (LLM) & C2 (Template) \\
\midrule
Contextual grounding    & 0.5667 & 0.4883 \\
Vocabulary diversity (TTR) & 0.3225 & 0.2766 \\
Unique words            & 396    & 151    \\
Avg.\ narration length  & 61.4 words & 27.3 words \\
Lexical coherence       & 0.2017 & 0.2127 \\
\bottomrule
\end{tabular}
\end{table}

LLM narration is more contextually grounded ($0.567$ vs.\ $0.488$), uses $2.6\times$ more unique words ($396$ vs.\ $151$), and produces $2.2\times$ longer descriptions ($61.4$ vs.\ $27.3$ words), while maintaining similar lexical coherence ($0.202$ vs.\ $0.213$). The slight coherence advantage for templates is expected: templates reuse fixed vocabulary, which mechanically inflates token-overlap similarity.
